% LaTeX Template for AI Problem Analysis Output (v2.1)
% AMC12A 2024 Problem 20 - Strategic Analysis
% Generated Analysis using AMC12/COMC Strategic Problem Analysis Framework

\documentclass[11pt]{article}
\usepackage[utf8]{inputenc}
\usepackage{amsmath, amssymb, amsthm}
\usepackage{geometry}
\usepackage{xcolor}
\usepackage[most]{tcolorbox}
\usepackage{enumitem}
\usepackage{fancyhdr}
\usepackage{hyperref}

% Page setup
\geometry{margin=1in}
\setlength{\headheight}{14.5pt}
\pagestyle{fancy}
\fancyhf{}
\rhead{AMC12 Problem Analysis}
\lhead{2024 AMC 12A Problem 20}
\cfoot{\thepage}

% Color definitions
\definecolor{problemconfig}{RGB}{230, 242, 255}    % Light blue
\definecolor{strategic}{RGB}{255, 230, 242}       % Light pink
\definecolor{tactical}{RGB}{242, 255, 230}        % Light green
\definecolor{techniques}{RGB}{255, 242, 230}      % Light yellow
\definecolor{verification}{RGB}{255, 230, 230}    % Light red
\definecolor{solution}{RGB}{230, 255, 255}        % Light cyan
\definecolor{competition}{RGB}{242, 242, 255}     % Light purple
\definecolor{gold}{RGB}{218, 165, 32}
\definecolor{highlight}{RGB}{255, 255, 200}       % Light yellow highlight

% Box styles
\newtcolorbox{problemconfigbox}{
    colback=problemconfig,
    colframe=blue,
    title=PROBLEM CONFIGURATION ANALYSIS,
    fonttitle=\bfseries,
    arc=3pt,
    boxrule=1pt,
    breakable
}

\newtcolorbox{strategicbox}{
    colback=strategic,
    colframe=purple,
    title=STRATEGIC FRAMEWORK APPLICATION,
    fonttitle=\bfseries,
    arc=3pt,
    boxrule=1pt,
    breakable
}

\newtcolorbox{tacticalbox}{
    colback=tactical,
    colframe=green,
    title=TACTICAL METHODOLOGY SELECTION,
    fonttitle=\bfseries,
    arc=3pt,
    boxrule=1pt,
    breakable
}

\newtcolorbox{techniquesbox}{
    colback=techniques,
    colframe=orange,
    title=CONVERSION TECHNIQUES APPLICATION,
    fonttitle=\bfseries,
    arc=3pt,
    boxrule=1pt,
    breakable
}

\newtcolorbox{verificationbox}{
    colback=verification,
    colframe=red,
    title=VERIFICATION STRATEGY DESIGN,
    fonttitle=\bfseries,
    arc=3pt,
    boxrule=1pt,
    breakable
}

\newtcolorbox{solutionbox}{
    colback=solution,
    colframe=cyan,
    title=SOLUTION ROADMAP,
    fonttitle=\bfseries,
    arc=3pt,
    boxrule=1pt,
    breakable
}

\newtcolorbox{competitionbox}{
    colback=competition,
    colframe=violet,
    title=COMPETITION STRATEGY INTEGRATION,
    fonttitle=\bfseries,
    arc=3pt,
    boxrule=1pt,
    breakable
}

\newtcolorbox{keyinsight}{
    colback=highlight,
    colframe=gold,
    title=KEY INSIGHT,
    fonttitle=\bfseries,
    arc=2pt,
    boxrule=2pt,
    breakable
}

\newtcolorbox{warningbox}{
    colback=red!10,
    colframe=red!80!black,
    title=WARNING,
    fonttitle=\bfseries,
    arc=2pt,
    boxrule=2pt,
    breakable
}

\newtcolorbox{protip}{
    colback=green!10,
    colframe=green!60!black,
    title=PRO TIP,
    fonttitle=\bfseries,
    arc=2pt,
    boxrule=2pt,
    breakable
}

\begin{document}

\title{\textbf{Strategic Analysis of AMC 12A 2024 Problem 20}}
\author{Using AMC12/COMC Strategic Problem Analysis Framework v2.1}
\date{}
\maketitle

\section*{Problem Statement}

Points $P$ and $Q$ are chosen uniformly and independently at random on sides $\overline{AB}$ and $\overline{AC}$, respectively, of equilateral triangle $\triangle ABC$. Which of the following intervals contains the probability that the area of $\triangle APQ$ is less than half the area of $\triangle ABC$?

\begin{align*}
\textbf{(A) } &\left[\frac{3}{8}, \frac{1}{2}\right] \qquad 
\textbf{(B) } \left(\frac{1}{2}, \frac{2}{3}\right] \qquad 
\textbf{(C) } \left(\frac{2}{3}, \frac{3}{4}\right] \\
\textbf{(D) } &\left(\frac{3}{4}, \frac{7}{8}\right] \qquad 
\textbf{(E) } \left(\frac{7}{8}, 1\right]
\end{align*}

\vspace{0.5cm}

\begin{problemconfigbox}
\subsection*{A. Problem Classification}
\begin{itemize}[leftmargin=*]
    \item \textbf{Problem Type}: To Find (computational/estimation problem)
    \item \textbf{Competition Context}: AMC 12A, Problem 20
    \item \textbf{Difficulty Band}: Late-band (P17-25), specifically P20 - one of the hardest problems
    \item \textbf{Mathematical Domain}: Probability, Geometry, Integral Calculus (optional)
    \item \textbf{Estimated Time}: 6-8 minutes for experienced students; 10-12 minutes for those learning the technique
\end{itemize}

\subsection*{B. Problem Structure Identification}
\begin{itemize}[leftmargin=*]
    \item \textbf{Given Information}: 
    \begin{itemize}
        \item Equilateral triangle $\triangle ABC$
        \item Point $P$ chosen uniformly at random on $\overline{AB}$
        \item Point $Q$ chosen uniformly at random on $\overline{AC}$
        \item $P$ and $Q$ are independent random variables
    \end{itemize}
    
    \item \textbf{Constraints}: 
    \begin{itemize}
        \item $P \in \overline{AB}$, $Q \in \overline{AC}$ (boundary constraints)
        \item Uniform distribution for both points (geometric probability)
        \item Independence of $P$ and $Q$ selections
    \end{itemize}
    
    \item \textbf{Objective}: Find which interval contains the probability that $[\triangle APQ] < \frac{1}{2}[\triangle ABC]$
    
    \item \textbf{Hidden Assumptions}: 
    \begin{itemize}
        \item The answer asks for an interval, not exact value (estimation acceptable)
        \item Uniform distribution on line segments corresponds to uniform distribution of distance parameters
        \item Area ratio depends on product of two independent uniform random variables
    \end{itemize}
    
    \item \textbf{Answer Choice Leverage}: 
    \begin{itemize}
        \item Intervals are well-separated, so rough estimation suffices
        \item Can use visual/graphical reasoning to eliminate choices
        \item Exact computation possible but not strictly necessary
        \item Pattern: intervals partition $[3/8, 1]$ with increasing widths
    \end{itemize}
\end{itemize}

\subsection*{C. Strategic Orientation}
\begin{itemize}[leftmargin=*]
    \item \textbf{Hypothesis}: If we parameterize $AP = x$ and $AQ = y$ (with appropriate normalization), the condition becomes an inequality in the unit square
    
    \item \textbf{Conclusion}: The probability equals the area of a region in the unit square $[0,1] \times [0,1]$
    
    \item \textbf{Notation}: 
    \begin{itemize}
        \item Let side length of $\triangle ABC$ be 1 (WLOG normalization)
        \item Let $AP = x$ where $x \in [0,1]$
        \item Let $AQ = y$ where $y \in [0,1]$
        \item $[\triangle ABC] = \frac{\sqrt{3}}{4}$ (area of unit equilateral triangle)
    \end{itemize}
    
    \item \textbf{Intuitive Approach}: 
    \begin{itemize}
        \item Express area of $\triangle APQ$ in terms of $x$ and $y$
        \item Use the formula: Area = $\frac{1}{2} \cdot AP \cdot AQ \cdot \sin(\angle PAQ)$
        \item Convert the area condition to an inequality: $xy \leq c$ for some constant $c$
        \item Visualize as a region in the unit square
        \item Estimate the area of this region
    \end{itemize}
    
    \item \textbf{Cross-Domain Potential}: 
    \begin{itemize}
        \item Geometric Probability: Converting to area computation in sample space
        \item Coordinate Geometry: Using $(x,y)$ parameterization
        \item Calculus: Computing area via integration (optional but more precise)
        \item Estimation/Visualization: Graphical area approximation
    \end{itemize}
\end{itemize}
\end{problemconfigbox}

\begin{keyinsight}
\textbf{The Core Insight}: This is a geometric probability problem that transforms into computing the area under a hyperbola in the unit square. The key realization is that the area condition $[\triangle APQ] < \frac{1}{2}[\triangle ABC]$ becomes the simple inequality $xy < \frac{1}{2}$ in normalized coordinates.
\end{keyinsight}

\begin{strategicbox}
\subsection*{A. Psychological Strategies}
\begin{itemize}[leftmargin=*]
    \item \textbf{Mental Toughness}: 
    \begin{itemize}
        \item This is P20 - expect it to be challenging
        \item Don't panic if the calculus approach seems difficult
        \item Visual/geometric reasoning can provide the answer without integration
        \item Trust that estimation is sufficient given the answer format
    \end{itemize}
    
    \item \textbf{Creativity Cultivation}: 
    \begin{itemize}
        \item Multiple approaches available: exact integration vs. visual estimation
        \item Can use symmetry and known geometric shapes to bound the area
        \item Consider what the graph of $xy = \frac{1}{2}$ looks like in the unit square
        \item Approximation by simpler curves (e.g., circular arc) can give bounds
    \end{itemize}
    
    \item \textbf{Iterative Refinement}: 
    \begin{itemize}
        \item Start with rough sketch to eliminate obviously wrong answers
        \item Refine estimate by comparing to known areas (half-square, quarter-circle)
        \item If time permits, compute exactly via integration
    \end{itemize}
\end{itemize}

\subsection*{B. Investigation Strategies}
\begin{itemize}[leftmargin=*]
    \item \textbf{Startup Orientation}: 
    \begin{itemize}
        \item Read carefully: ``uniformly and independently at random''
        \item Identify that this is a 2D geometric probability problem
        \item Recognize the need to parameterize positions of $P$ and $Q$
    \end{itemize}
    
    \item \textbf{Get Your Hands Dirty}: 
    \begin{itemize}
        \item Draw an equilateral triangle with points $P$ and $Q$
        \item Try specific positions: $P$ at midpoint, $Q$ at midpoint
        \item Calculate area for this specific case
        \item Ask: when is the area exactly half? This gives the boundary curve
    \end{itemize}
    
    \item \textbf{Penultimate Step}: 
    \begin{itemize}
        \item The final step is comparing our computed/estimated probability to the given intervals
        \item The penultimate step is determining the area of region $\{(x,y): xy \leq \frac{1}{2}, x,y \in [0,1]\}$
    \end{itemize}
    
    \item \textbf{Make It Easier}: 
    \begin{itemize}
        \item Normalize the triangle to have side length 1
        \item Use the fact that angles in equilateral triangle are known ($60^\circ$)
        \item Recognize that sine formula gives simple expression
    \end{itemize}
    
    \item \textbf{Conjecture via Small Cases}: 
    \begin{itemize}
        \item When $P = A$ or $Q = A$: area is 0 (always less than half)
        \item When $P = B$ and $Q = C$: area equals full triangle (not less than half)
        \item When $P$ at midpoint and $Q$ at midpoint: $xy = \frac{1}{4}$, area = $\frac{1}{4}$ of original (less than half)
    \end{itemize}
    
    \item \textbf{Visualization}: 
    \begin{itemize}
        \item Draw the unit square $[0,1] \times [0,1]$ representing $(x,y)$ space
        \item Sketch the curve $xy = \frac{1}{2}$, which is a hyperbola
        \item The curve passes through $(0.5, 1)$ and $(1, 0.5)$
        \item Shade the region where $xy \leq \frac{1}{2}$
        \item Compare visually to half-square ($0.5$) and $\frac{3}{4}$ square
    \end{itemize}
\end{itemize}

\subsection*{C. Argument Construction (Late-Band Playbook)}
\begin{itemize}[leftmargin=*]
    \item \textbf{Late-Band Recognition}: This is P20, solidly in the late band
    
    \item \textbf{Pattern Matching}: 
    \begin{itemize}
        \item Geometric probability + uniform distribution → area computation
        \item Triangle area with two sides and included angle → sine formula
        \item Inequality constraint → region in parameter space
    \end{itemize}
    
    \item \textbf{Transformation Chain}: 
    \begin{itemize}
        \item Physical setup → Parameterization ($x$, $y$)
        \item Area condition → Algebraic inequality ($xy \leq \frac{1}{2}$)
        \item Probability → Geometric area in unit square
        \item Exact or estimated area → Answer interval selection
    \end{itemize}
    
    \item \textbf{Answer Choice Strategy}: 
    \begin{itemize}
        \item The intervals are mutually exclusive and cover different ranges
        \item Visual inspection can eliminate (A), (B), (E)
        \item Distinguish between (C) and (D) requires better estimate
        \item Integration gives $\frac{\ln 2 + 1}{2} \approx 0.8466$ → answer (D)
    \end{itemize}
\end{itemize}
\end{strategicbox}

\begin{warningbox}
\textbf{Common Pitfall \#1}: Forgetting the independence assumption. The probability is NOT the product of individual probabilities, but rather the area of a region in the joint sample space.

\textbf{Common Pitfall \#2}: Incorrect area formula. Must use $\frac{1}{2}xy\sin(60^\circ) = \frac{\sqrt{3}}{4}xy$, not just $\frac{1}{2}xy$.

\textbf{Common Pitfall \#3}: Computing exact integral when rough estimation suffices. This wastes time on a test.
\end{warningbox}

\begin{tacticalbox}
\subsection*{A. Core Mathematical Tactics}
\begin{itemize}[leftmargin=*]
    \item \textbf{Normalization Principle}: 
    \begin{itemize}
        \item Set side length to 1 (without loss of generality)
        \item This simplifies all calculations
        \item Ratios and probabilities are preserved under scaling
    \end{itemize}
    
    \item \textbf{Symmetry Exploitation}: 
    \begin{itemize}
        \item The curve $xy = \frac{1}{2}$ has symmetry about the line $y = x$
        \item This doesn't directly help but confirms the aesthetic of the problem
    \end{itemize}
    
    \item \textbf{Extreme Case Analysis}: 
    \begin{itemize}
        \item When $x = 0$ or $y = 0$: always satisfied ($xy = 0 < \frac{1}{2}$)
        \item When $x = 1, y = 1$: $xy = 1 > \frac{1}{2}$ (not satisfied)
        \item This confirms we're looking at a proper subset of the unit square
    \end{itemize}
    
    \item \textbf{Dimensional Analysis}: 
    \begin{itemize}
        \item Areas scale as length squared
        \item Probability is dimensionless (ratio of areas)
        \item This confirms our approach is dimensionally consistent
    \end{itemize}
    
    \item \textbf{Invariance Recognition}: 
    \begin{itemize}
        \item The probability doesn't depend on which two sides we choose
        \item Any two sides meeting at a vertex would give the same answer
    \end{itemize}
\end{itemize}

\subsection*{B. Crossover Techniques}
\begin{itemize}[leftmargin=*]
    \item \textbf{Geometry → Algebra}: 
    \begin{itemize}
        \item Geometric constraint (area inequality) becomes algebraic inequality ($xy \leq \frac{1}{2}$)
    \end{itemize}
    
    \item \textbf{Probability → Geometry}: 
    \begin{itemize}
        \item Uniform probability on segments maps to uniform distribution on unit square
        \item Probability computation reduces to area computation
    \end{itemize}
    
    \item \textbf{Continuous → Discrete}: 
    \begin{itemize}
        \item While the problem involves continuous distributions, the answer is discrete (interval selection)
    \end{itemize}
    
    \item \textbf{Exact → Approximate}: 
    \begin{itemize}
        \item Can solve exactly with calculus
        \item Can estimate visually for interval identification
        \item Strategy choice depends on time constraints
    \end{itemize}
\end{itemize}
\end{tacticalbox}

\begin{techniquesbox}
\subsection*{A. Recognition Index}
\begin{itemize}[leftmargin=*]
    \item \textbf{Primary Technique}: Geometric Probability via Area Computation
    \item \textbf{Secondary Technique}: Parametric Representation of Triangle Area
    \item \textbf{Advanced Technique}: Integration under a curve (optional)
    \item \textbf{Estimation Technique}: Visual approximation and bounds
\end{itemize}

\subsection*{B. Technique Selection}
\begin{itemize}[leftmargin=*]
    \item \textbf{Step 1 - Parameterization}: 
    \begin{itemize}
        \item Represent $P$ by parameter $x = \frac{AP}{AB}$ where $x \in [0,1]$
        \item Represent $Q$ by parameter $y = \frac{AQ}{AC}$ where $y \in [0,1]$
        \item Sample space is the unit square $[0,1] \times [0,1]$
    \end{itemize}
    
    \item \textbf{Step 2 - Area Formula}: 
    \begin{itemize}
        \item Use sine formula: $[\triangle APQ] = \frac{1}{2} \cdot AP \cdot AQ \cdot \sin(\angle A)$
        \item For equilateral triangle: $\angle A = 60^\circ$, $\sin(60^\circ) = \frac{\sqrt{3}}{2}$
        \item With normalization $AB = AC = 1$: $AP = x$, $AQ = y$
        \item Therefore: $[\triangle APQ] = \frac{1}{2}xy \cdot \frac{\sqrt{3}}{2} = \frac{\sqrt{3}}{4}xy$
    \end{itemize}
    
    \item \textbf{Step 3 - Inequality Conversion}: 
    \begin{itemize}
        \item Given: $[\triangle ABC] = \frac{\sqrt{3}}{4}$ (unit equilateral triangle)
        \item Condition: $[\triangle APQ] < \frac{1}{2}[\triangle ABC]$
        \item Substituting: $\frac{\sqrt{3}}{4}xy < \frac{1}{2} \cdot \frac{\sqrt{3}}{4}$
        \item Simplifying: $xy < \frac{1}{2}$
    \end{itemize}
    
    \item \textbf{Step 4 - Geometric Probability}: 
    \begin{itemize}
        \item Probability = $\frac{\text{Area of favorable region}}{\text{Area of sample space}}$
        \item Sample space area = $1 \times 1 = 1$
        \item Favorable region: $\{(x,y): 0 \leq x \leq 1, 0 \leq y \leq 1, xy \leq \frac{1}{2}\}$
    \end{itemize}
    
    \item \textbf{Step 5 - Area Computation (Method A - Integration)}: 
    \begin{itemize}
        \item The boundary curve is $y = \frac{1}{2x}$
        \item For $0 < x < \frac{1}{2}$: $y$ ranges from 0 to 1 (curve above the square)
        \item For $\frac{1}{2} \leq x \leq 1$: $y$ ranges from 0 to $\frac{1}{2x}$ (curve inside square)
        \item Area = $\int_0^{1/2} 1 \, dx + \int_{1/2}^1 \frac{1}{2x} \, dx$
        \item $= \frac{1}{2} + \frac{1}{2}[\ln x]_{1/2}^1 = \frac{1}{2} + \frac{1}{2}(\ln 1 - \ln \frac{1}{2})$
        \item $= \frac{1}{2} + \frac{1}{2}(0 - (-\ln 2)) = \frac{1}{2} + \frac{\ln 2}{2} = \frac{1 + \ln 2}{2}$
        \item Numerical value: $\frac{1 + 0.693}{2} \approx 0.8466$
    \end{itemize}
    
    \item \textbf{Step 5 - Area Computation (Method B - Visual Estimation)}: 
    \begin{itemize}
        \item Lower bound: The region includes the entire left half-square and bottom half-square minus overlap
        \item This gives at least $\frac{1}{2} + \frac{1}{2} - \frac{1}{4} = \frac{3}{4}$
        \item Upper bound: The region is less than the whole square (1)
        \item The hyperbola $xy = \frac{1}{2}$ curves gently, bulging outward
        \item Visual inspection suggests area is close to $\frac{7}{8}$ but slightly less
        \item This places the answer in interval (D): $(\frac{3}{4}, \frac{7}{8}]$
    \end{itemize}
\end{itemize}

\subsection*{C. Integration \& Conflict Resolution}
\begin{itemize}[leftmargin=*]
    \item \textbf{Consistency Check}: 
    \begin{itemize}
        \item Exact value $\frac{1 + \ln 2}{2} \approx 0.8466$ confirms interval (D)
        \item Visual estimate consistent with exact calculation
        \item No conflicts between approaches
    \end{itemize}
    
    \item \textbf{Trade-offs}: 
    \begin{itemize}
        \item Integration gives exact answer but requires calculus and time
        \item Visual estimation is faster but requires good spatial intuition
        \item On an exam, visual approach adequate for interval selection
    \end{itemize}
\end{itemize}
\end{techniquesbox}

\begin{protip}
\textbf{Time-Saving Strategy}: On exam day, sketch the curve $y = \frac{1}{2x}$ in the unit square. Observe it passes through $(0.5, 1)$ and $(1, 0.5)$. The region below the curve is clearly more than $\frac{3}{4}$ of the square (eliminates A, B, C) but less than $\frac{7}{8}$ (eliminates E). Select (D) without computing the integral. \textbf{Estimated time savings: 3-4 minutes.}
\end{protip}

\begin{verificationbox}
\subsection*{A. Verification Level Selection}
\begin{itemize}[leftmargin=*]
    \item \textbf{Chosen Level}: Level 4 (Sanity Check + Alternative Method)
    \item \textbf{Rationale}: This is P20 (late band, high stakes), and we need high confidence
    \item \textbf{Time Allocation}: 1-2 minutes
\end{itemize}

\subsection*{B. Verification Methods}

\textbf{Method 1: Boundary Check}
\begin{itemize}
    \item At $(0.5, 1)$: $xy = 0.5$, exactly on boundary $\checkmark$
    \item At $(1, 0.5)$: $xy = 0.5$, exactly on boundary $\checkmark$
    \item At $(1, 1)$: $xy = 1 > 0.5$, correctly excluded $\checkmark$
    \item At $(0, y)$ or $(x, 0)$: $xy = 0 < 0.5$, correctly included $\checkmark$
\end{itemize}

\textbf{Method 2: Dimensional Analysis}
\begin{itemize}
    \item Probability must be dimensionless $\checkmark$
    \item Probability must be in $[0, 1]$ $\checkmark$
    \item Our answer $\approx 0.8466 \in [0,1]$ $\checkmark$
\end{itemize}

\textbf{Method 3: Limiting Cases}
\begin{itemize}
    \item If we required $[\triangle APQ] < \varepsilon$ (very small): probability $\to 0$ $\checkmark$
    \item If we required $[\triangle APQ] < [\triangle ABC]$: probability $\to 1$ $\checkmark$
    \item Our condition is intermediate, so probability between 0 and 1 is reasonable $\checkmark$
\end{itemize}

\textbf{Method 4: Symmetry Verification}
\begin{itemize}
    \item The curve $xy = c$ is symmetric about $y = x$ $\checkmark$
    \item Our computation doesn't explicitly use this, but it's consistent
    \item If we chose $P$ on $AC$ and $Q$ on $AB$ instead, we'd get the same answer by symmetry $\checkmark$
\end{itemize}

\textbf{Method 5: Answer Choice Elimination}
\begin{itemize}
    \item Visual inspection clearly shows area $> 0.75$ (eliminates A, B, C)
    \item Visual inspection suggests area $< 0.875$ (eliminates E)
    \item Only (D) remains
    \item Exact calculation confirms (D)
\end{itemize}

\subsection*{C. Confidence Calibration}
\begin{itemize}[leftmargin=*]
    \item \textbf{Initial Confidence}: High (95\%) - standard geometric probability problem
    \item \textbf{After Verification}: Very High (99\%) - multiple methods agree
    \item \textbf{Final Decision}: Submit answer (D) with high confidence
\end{itemize}
\end{verificationbox}

\begin{solutionbox}
\subsection*{Step-by-Step Solution Plan}

\textbf{Phase 1: Setup and Parameterization (1-2 minutes)}
\begin{enumerate}
    \item Draw equilateral triangle $\triangle ABC$
    \item Mark points $P$ on $\overline{AB}$ and $Q$ on $\overline{AC}$
    \item Set side length to 1 (WLOG)
    \item Define parameters: $x = AP$, $y = AQ$ where $x, y \in [0,1]$
    \item Recognize: uniform random selection $\Rightarrow$ uniform distribution on $[0,1] \times [0,1]$
\end{enumerate}

\textbf{Phase 2: Area Formula Derivation (1-2 minutes)}
\begin{enumerate}[resume]
    \item Apply sine area formula: $[\triangle APQ] = \frac{1}{2} \cdot AP \cdot AQ \cdot \sin(\angle PAQ)$
    \item Note: $\angle PAQ = 60^\circ$ (angle of equilateral triangle)
    \item Substitute: $[\triangle APQ] = \frac{1}{2} \cdot x \cdot y \cdot \frac{\sqrt{3}}{2} = \frac{\sqrt{3}}{4}xy$
    \item Calculate: $[\triangle ABC] = \frac{\sqrt{3}}{4}$ (standard formula for unit equilateral triangle)
\end{enumerate}

\textbf{Phase 3: Inequality Conversion (1 minute)}
\begin{enumerate}[resume]
    \item Set up condition: $[\triangle APQ] < \frac{1}{2}[\triangle ABC]$
    \item Substitute: $\frac{\sqrt{3}}{4}xy < \frac{1}{2} \cdot \frac{\sqrt{3}}{4}$
    \item Simplify: $xy < \frac{1}{2}$
    \item \textbf{Key insight}: The geometric condition reduces to a simple algebraic inequality
\end{enumerate}

\textbf{Phase 4: Probability Computation (2-4 minutes)}

\textit{Option A: Visual Estimation (faster, sufficient for exam)}
\begin{enumerate}[resume]
    \item Sketch unit square with axes $x$ and $y$
    \item Draw hyperbola $xy = \frac{1}{2}$ passing through $(0.5, 1)$ and $(1, 0.5)$
    \item Shade region where $xy \leq \frac{1}{2}$ (below/left of curve)
    \item Observe: region includes entire left half and bottom half with overlap
    \item Estimate visually: area appears to be approximately 0.85
    \item Compare to intervals: clearly in $(\frac{3}{4}, \frac{7}{8}]$
\end{enumerate}

\textit{Option B: Exact Integration (more rigorous)}
\begin{enumerate}[resume]
    \item Set up integral: Area $= \int_0^{1/2} 1 \, dx + \int_{1/2}^1 \frac{1}{2x} \, dx$
    \item Compute first integral: $\int_0^{1/2} 1 \, dx = \frac{1}{2}$
    \item Compute second integral: $\int_{1/2}^1 \frac{1}{2x} \, dx = \frac{1}{2}[\ln x]_{1/2}^1 = \frac{1}{2}(0 - \ln\frac{1}{2}) = \frac{\ln 2}{2}$
    \item Sum: Total area $= \frac{1}{2} + \frac{\ln 2}{2} = \frac{1 + \ln 2}{2}$
    \item Evaluate: Using $\ln 2 \approx 0.693$, we get $\frac{1 + 0.693}{2} \approx 0.8466$
    \item Identify interval: $0.75 < 0.8466 < 0.875$, so answer is (D)
\end{enumerate}

\textbf{Phase 5: Verification (1-2 minutes)}
\begin{enumerate}[resume]
    \item Check boundary points: $(0.5, 1)$ and $(1, 0.5)$ should satisfy $xy = 0.5$ $\checkmark$
    \item Verify probability is in $[0,1]$: $0.8466 \in [0,1]$ $\checkmark$
    \item Check reasonableness: Most of the unit square should satisfy $xy < 0.5$ $\checkmark$
    \item Confirm answer choice: (D) $\left(\frac{3}{4}, \frac{7}{8}\right]$ $\checkmark$
\end{enumerate}

\subsection*{Alternative Approaches}

\textbf{Alternative 1: Monte Carlo Simulation (if calculator allowed)}
\begin{itemize}
    \item Generate many random pairs $(x, y)$ from $[0,1] \times [0,1]$
    \item Count proportion where $xy < 0.5$
    \item Should converge to approximately 0.846
\end{itemize}

\textbf{Alternative 2: Complementary Probability}
\begin{itemize}
    \item Compute $P(xy > \frac{1}{2})$ instead
    \item This is the area above the hyperbola: $\int_0^{1/2} 0 \, dx + \int_{1/2}^1 (1 - \frac{1}{2x}) \, dx$
    \item $= \int_{1/2}^1 (1 - \frac{1}{2x}) \, dx = [x - \frac{1}{2}\ln x]_{1/2}^1$
    \item $= (1 - 0) - (\frac{1}{2} - \frac{1}{2}\ln\frac{1}{2}) = 1 - \frac{1}{2} - \frac{\ln 2}{2} = \frac{1 - \ln 2}{2}$
    \item Then $P(xy \leq \frac{1}{2}) = 1 - \frac{1 - \ln 2}{2} = \frac{1 + \ln 2}{2}$ $\checkmark$
\end{itemize}

\subsection*{Common Pitfalls and How to Avoid Them}

\begin{enumerate}
    \item \textbf{Pitfall}: Using $[\triangle APQ] = \frac{1}{2}xy$ without the $\sin(60^\circ)$ factor
    \begin{itemize}
        \item \textbf{Consequence}: Wrong inequality $xy < \frac{\sqrt{3}}{2}$
        \item \textbf{Prevention}: Always write out the sine formula explicitly
    \end{itemize}
    
    \item \textbf{Pitfall}: Forgetting to normalize the triangle side length
    \begin{itemize}
        \item \textbf{Consequence}: Messy algebra with arbitrary constants
        \item \textbf{Prevention}: State "WLOG, let side length = 1" at the beginning
    \end{itemize}
    
    \item \textbf{Pitfall}: Thinking the probability is $P(x < c) \cdot P(y < c')$ for some constants
    \begin{itemize}
        \item \textbf{Consequence}: Wrong model of the problem
        \item \textbf{Prevention}: Recognize this is about the joint distribution, not independent events
    \end{itemize}
    
    \item \textbf{Pitfall}: Computational error in the integral
    \begin{itemize}
        \item \textbf{Consequence}: Wrong numerical value
        \item \textbf{Prevention}: Double-check logarithm rules: $\ln(1/2) = -\ln 2$
    \end{itemize}
    
    \item \textbf{Pitfall}: Spending too much time on exact calculation
    \begin{itemize}
        \item \textbf{Consequence}: Time pressure on remaining problems
        \item \textbf{Prevention}: For P20, visual estimation is acceptable given interval format
    \end{itemize}
\end{enumerate}
\end{solutionbox}

\begin{competitionbox}
\subsection*{A. Time Management Strategy}
\begin{itemize}[leftmargin=*]
    \item \textbf{Target Time}: 6-8 minutes for this problem
    \item \textbf{Time Checkpoints}:
    \begin{itemize}
        \item 2 min: Should have parameterization and area formula
        \item 4 min: Should have inequality $xy < \frac{1}{2}$ and sketch drawn
        \item 6 min: Should have answer selected (D) via estimation or integration
        \item 8 min: Should complete verification and move on
    \end{itemize}
    \item \textbf{Adjustment Rule}: If stuck after 10 minutes, mark answer (D) based on visual estimate and move on
\end{itemize}

\subsection*{B. Prioritization}
\begin{itemize}[leftmargin=*]
    \item \textbf{Problem Positioning}: P20 is very late in the test
    \item \textbf{When to Attempt}: 
    \begin{itemize}
        \item Only after completing all problems you're confident about (P1-16)
        \item Only if you have at least 10 minutes remaining
        \item Skip if time is critically short
    \end{itemize}
    \item \textbf{Expected Success Rate}: 
    \begin{itemize}
        \item Target audience: Top 5-10\% of test-takers
        \item Requires: Geometric probability + integration OR strong spatial reasoning
    \end{itemize}
\end{itemize}

\subsection*{C. Risk Management}
\begin{itemize}[leftmargin=*]
    \item \textbf{Risk Level}: Medium-High (P20, but tractable with right approach)
    \item \textbf{Mitigation Strategies}:
    \begin{itemize}
        \item Use visual estimation instead of exact integration to save time
        \item Answer choice spacing is generous - rough estimate suffices
        \item Can eliminate obviously wrong intervals quickly
    \end{itemize}
    \item \textbf{Confidence Threshold}: Proceed if confidence $\geq$ 60\% after 4 minutes of work
\end{itemize}

\subsection*{D. Abandonment Criteria}
\begin{itemize}[leftmargin=*]
    \item \textbf{When to Abandon}:
    \begin{itemize}
        \item If unable to derive $xy < \frac{1}{2}$ after 5 minutes
        \item If cannot sketch or visualize the region after 6 minutes
        \item If less than 8 minutes remain in the exam
    \end{itemize}
    \item \textbf{Abandonment Strategy}:
    \begin{itemize}
        \item Make educated guess based on partial work
        \item Visual estimation suggests (D) is most likely even without full solution
        \item Return to problem only if time permits after checking all other answers
    \end{itemize}
\end{itemize}

\subsection*{E. Answer Recording}
\begin{itemize}[leftmargin=*]
    \item \textbf{Format Check}: Multiple choice - bubble in (D)
    \item \textbf{Double-Check}: Verify you've selected an interval, not a point value
    \item \textbf{Final Confirmation}: $\left(\frac{3}{4}, \frac{7}{8}\right]$ contains approximately 0.8466
\end{itemize}
\end{competitionbox}

\section*{Complete Worked Solution}

\subsection*{Approach 1: Visual Estimation (Exam Strategy - Fast)}

\textbf{Setup}: Let the equilateral triangle have side length 1 (WLOG). Let $P$ be at distance $x$ from $A$ along $\overline{AB}$, and $Q$ be at distance $y$ from $A$ along $\overline{AC}$, where $x, y \in [0,1]$.

\textbf{Area Formula}: Using the formula for the area of a triangle with two sides and included angle:
\[[\triangle APQ] = \frac{1}{2} \cdot x \cdot y \cdot \sin(60^\circ) = \frac{1}{2} \cdot x \cdot y \cdot \frac{\sqrt{3}}{2} = \frac{\sqrt{3}}{4}xy\]

Since $[\triangle ABC] = \frac{\sqrt{3}}{4}$, the condition $[\triangle APQ] < \frac{1}{2}[\triangle ABC]$ becomes:
\[\frac{\sqrt{3}}{4}xy < \frac{1}{2} \cdot \frac{\sqrt{3}}{4}\]
\[xy < \frac{1}{2}\]

\textbf{Probability as Area}: The sample space is the unit square $[0,1] \times [0,1]$ with area 1. The favorable region is $\{(x,y): xy \leq \frac{1}{2}\}$.

\textbf{Visual Analysis}: The boundary curve $xy = \frac{1}{2}$ is a hyperbola passing through $(0.5, 1)$ and $(1, 0.5)$. Sketching this curve:
\begin{itemize}
    \item The region below/left of the curve clearly includes more than $\frac{3}{4}$ of the square
    \item But it's less than $\frac{7}{8}$ of the square (the hyperbola cuts off a noticeable corner)
    \item Visual estimation: approximately 0.84-0.85 of the square
\end{itemize}

\textbf{Answer}: The interval $\left(\frac{3}{4}, \frac{7}{8}\right]$ contains this probability. $\boxed{\textbf{(D)}}$

\subsection*{Approach 2: Exact Integration (Rigorous)}

\textbf{Setup}: Same as Approach 1, leading to $xy < \frac{1}{2}$.

\textbf{Integration Setup}: The favorable region can be computed as:
\[\text{Area} = \int_0^1 \min\left(1, \frac{1}{2x}\right) dx\]

Split the integral at $x = \frac{1}{2}$ (where the curve intersects $y = 1$):
\begin{align*}
\text{Area} &= \int_0^{1/2} 1 \, dx + \int_{1/2}^1 \frac{1}{2x} \, dx \\
&= \left[x\right]_0^{1/2} + \frac{1}{2}\left[\ln x\right]_{1/2}^1 \\
&= \frac{1}{2} + \frac{1}{2}(\ln 1 - \ln \frac{1}{2}) \\
&= \frac{1}{2} + \frac{1}{2}(0 - (-\ln 2)) \\
&= \frac{1}{2} + \frac{\ln 2}{2} \\
&= \frac{1 + \ln 2}{2}
\end{align*}

\textbf{Numerical Evaluation}: Using $\ln 2 \approx 0.693147$:
\[\frac{1 + \ln 2}{2} \approx \frac{1 + 0.693}{2} \approx \frac{1.693}{2} \approx 0.8466\]

\textbf{Interval Check}: 
\[\frac{3}{4} = 0.75 < 0.8466 < 0.875 = \frac{7}{8}\]

Therefore, the probability lies in the interval $\left(\frac{3}{4}, \frac{7}{8}\right]$.

\textbf{Answer}: $\boxed{\textbf{(D) } \left(\frac{3}{4}, \frac{7}{8}\right]}$

\subsection*{Approximation Note (for non-calculator exam)}

If you don't have $\ln 2$ memorized, you can use:
\[\ln 2 = \ln(1+1) \approx 1 - \frac{1}{2} + \frac{1}{3} - \frac{1}{4} = \frac{12-6+4-3}{12} = \frac{7}{12} \approx 0.583\]

Or more accurately (4 terms):
\[\ln 2 \approx 1 - \frac{1}{2} + \frac{1}{3} - \frac{1}{4} + \frac{1}{5} = \frac{60-30+20-15+12}{60} = \frac{47}{60} \approx 0.783\]

Wait, this series is for $\ln(1+x)$ with $x=1$, which doesn't converge well. Better: use $\ln 2 \approx 0.69$ from memory or note that the answer must be in (D) from visual estimation alone.

\section*{Key Takeaways}

\begin{enumerate}
    \item \textbf{Geometric probability} often reduces to computing areas in parameter space
    \item \textbf{Parameterization} is key: representing positions by normalized distances
    \item \textbf{The sine formula} for triangle area is essential when two sides and included angle are known
    \item \textbf{Visual estimation} can be sufficient for interval-based answer choices
    \item \textbf{Integration} provides exact answers but takes more time
    \item \textbf{Late-band problems} (P17-25) often require multiple steps and sophisticated reasoning
    \item \textbf{Answer intervals} being well-separated means rough estimates are acceptable
\end{enumerate}

\section*{Final Answer}

The probability that the area of $\triangle APQ$ is less than half the area of $\triangle ABC$ is exactly $\dfrac{1 + \ln 2}{2} \approx 0.8466$, which lies in the interval $\left(\dfrac{3}{4}, \dfrac{7}{8}\right]$.

\begin{center}
\fbox{\Large \textbf{Answer: (D) } $\left(\dfrac{3}{4}, \dfrac{7}{8}\right]$}
\end{center}

\end{document}